% Example "decay": the problem, the solutions against the exact one, and the
% convergence test. Every number comes from data/ (parameters.tex,
% results.tex, exact.dat, solution_<method>.dat, error_<method>.dat), written
% by examples/decay/main.f90: run `fpm run --example decay` from the project
% root, then build this file in examples/decay/.
\documentclass[11pt,a4paper]{article}
\usepackage[margin=2.2cm]{geometry}
\usepackage{amsmath,booktabs}
\usepackage{pgfplots}
\pgfplotsset{compat=1.18}
\pagestyle{empty}

\begin{document}
\section*{Example: exponential decay}

\subsection*{1. Problem}
\begin{equation*}
  \frac{dy}{dt} = -\lambda\, y, \qquad y(0) = 1,
  \qquad\text{exact solution}\qquad y(t) = e^{-\lambda t}.
\end{equation*}
The simplest linear problem: every Runge--Kutta step multiplies $y$ by the
stability polynomial of the method evaluated at $z = -\lambda\Delta t$, so the
numerical solution is $y_n = R(-\lambda\Delta t)^n$ and the whole error comes
from how well $R(z)$ approximates $e^{z}$.
\begin{center}
% Written by examples/decay/main.f90: do not edit.
\begin{tabular}{ll}
\toprule
decay rate & $\lambda = 2.00$ \\
initial condition & $y(0) = 1$ \\
interval & $0 \le t \le 1.00$ \\
steps, convergence & $\Delta t = 1/10, 1/20, \dots, 1/1280$ (8 levels) \\
step, solution figure & $\Delta t = 0.25$ (4 steps) \\
methods & \texttt{rk3ws}, \texttt{rk4}, \texttt{dp54} \\
interface & model classes, \texttt{rk\_scheme(scheme, name, dt, model)} \\
observed order & slope of the two finest levels with errors above $10^{-12}$ \\
\bottomrule
\end{tabular}

\end{center}

\subsection*{2. Solutions}
\begin{center}
\begin{tikzpicture}
  \begin{axis}[
      width=11cm, height=6.5cm,                            % figure size
      xlabel={$t$}, ylabel={$y$},                          % axes: time, solution
      xmin=0, xmax=1, ymin=0, ymax=1.05,                   % the whole interval
      grid=major, grid style={black!12},                   % light grid
      legend pos=outer north east, legend cell align=left, % legend right of the axes
      legend style={font=\small},                          % compact legend
    ]
    \addplot[black, thick] table[x=t, y=y] {data/exact.dat};              % exact e^{-lambda t}
    \addlegendentry{exact}
    \addplot[green!50!black, mark=triangle*, mark size=1.8pt] table[x=t, y=y] {data/solution_rk3ws.dat};  % order 3
    \addlegendentry{\texttt{rk3ws}}
    \addplot[blue!80!black, mark=diamond*, mark size=1.8pt] table[x=t, y=y] {data/solution_rk4.dat};       % order 4
    \addlegendentry{\texttt{rk4}}
    \addplot[red!75!black, only marks, mark=o, mark size=2.5pt] table[x=t, y=y] {data/solution_dp54.dat}; % order 5
    \addlegendentry{\texttt{dp54}}
  \end{axis}
\end{tikzpicture}
\end{center}
At the coarse step of the figure the three methods of the library
(\texttt{rk3ws} is of order 3 on this linear problem) fall on the exact curve
at this scale; forward Euler, given as a tableau, would halve $y$ at each step
($R(-0.5) = 0.5$ against $e^{-0.5} \approx 0.61$).

\newpage
\subsection*{3. Test: convergence order}

\paragraph{Prediction.} A method of order $p$ has a global error
$E \approx C\,\Delta t^{p}$: halving $\Delta t$ divides it by $2^p$, a line of
slope $p$ in log--log axes. On a linear problem the order that counts is how
many terms of $e^{z}$ the polynomial $R(z)$ matches
($b^{\mathsf T}A^{k-1}\mathbf{1} = 1/k!$), which can exceed the general order:
\texttt{rk3ws} is of order 2 in general but 3 here.

\paragraph{Result.}
\begin{center}
% Written by examples/decay/main.f90: do not edit.
\begin{tabular}{llccc}
\toprule
method & name & order & linear order & observed \\
\midrule
RK3 of Wicker and Skamarock & \texttt{rk3ws} & 2 & 3 &  3.00 \\
Classic RK4 & \texttt{rk4} & 4 & 4 &  4.01 \\
Dormand-Prince 5(4) & \texttt{dp54} & 5 & 5 &  5.12 \\
\bottomrule
\end{tabular}

\end{center}

\begin{center}
\begin{tikzpicture}
  \begin{loglogaxis}[
      width=10.5cm, height=7.5cm,                          % figure size
      xlabel={time step $\Delta t$},                       % x: step
      ylabel={error $|y_N - e^{-\lambda}|$ at $t = 1$},    % y: global error
      grid=major, grid style={black!12},                   % light grid
      legend pos=outer north east, legend cell align=left, % legend right of the axes
      legend style={font=\small},                          % compact legend
      every axis plot/.append style={thick, mark size=1.5pt}, % same weight for all curves
    ]
    \addplot[green!50!black, dashed, mark=triangle] table[x=dt, y=error] {data/error_rk3ws.dat}; % rk3ws: slope 3 (linear)
    \addlegendentry{\texttt{rk3ws}}
    \addplot[blue!80!black, mark=diamond*]   table[x=dt, y=error] {data/error_rk4.dat};      % RK4: slope 4
    \addlegendentry{\texttt{rk4}}
    \addplot[red!75!black, mark=pentagon*]   table[x=dt, y=error] {data/error_dp54.dat};     % DP5(4): slope 5
    \addlegendentry{\texttt{dp54}}
  \end{loglogaxis}
\end{tikzpicture}
\end{center}
Every method converges at its predicted order. Curves stop where the error
reaches round-off; methods sharing a stability polynomial coincide.
\end{document}
